\section{Discussion}
\label{sec:limits}
R20 produced six complete task-to-approval-to-receipt-to-sample traces across
two simulator scenes; none of its single post-readback samples reached target.
R21 passed 31 assertions on the same configured route, and a separate mock
test refused execution after the approved active run changed. Together, these
results verify selected software behavior, not manufacturing performance.

The runs do not establish independent replication, process convergence, or
quality improvement. Each scene used one simulator instance, and its single
probe is not a response trajectory. Configuration reuse, alternative
harnesses and direct REST entry, normal online DAQ, operator effectiveness,
and physical transfer remain untested. Synthetic approvals do not assess
operator effectiveness \cite{parasuramanmanzey2010complacency,parasuraman2000model,bainbridge1983ironies}; restricted tool access does not establish deployment security \cite{greshake2023injection,liu2024pinjection,iec62443-3-3} or replace independently engineered PLC/SIS protection \cite{iec61511}.

Further work should measure configuration changes against a baseline and test
alternate routes, independent simulator runs, and the online DAQ backend.
Physical-plant trials are required to assess process outcomes.

\section{Conclusion}
\label{sec:conclusion}
AgentWorkShop connects collaborative tasks to production context, checked
operations, semantic I/O, and service-owned records. R20 and R21 verify
selected paths across two simulator mappings; a regression test confirms that
a stale approval cannot cross an active-run change. The evidence supports the
implemented software contract, not process-control efficacy, comparative
advantage, or portability.
